\begin{proof}
On the unit circle, $|P^*|=|P|$. Hence
\eqref{I:B:eq:curvature-hardy-residual} gives
\begin{equation}
 |d|^2=\frac{|P|^2}{C^2}\quad\text{almost everywhere}.
 \label{I:B:eq:curvature-boundary-modulus}
\end{equation}
Any boundary zeros of $P$ cancel in the rational expression for $d$;
they do not alter this identity or its coefficient consequences.
Integrating \eqref{I:B:eq:curvature-boundary-modulus} proves $A=C^{-2}$.

The $i$th scalar coefficient of $H_dS_+\overline d$ is
\begin{equation}
 (H_dS_+\overline d)_i
   =\sum_{j\ge0}d_{i+j+1}\overline{d_j},\qquad i\ge0.
 \label{I:B:eq:curvature-autocorrelation-d}
\end{equation}
These sums converge absolutely by Cauchy--Schwarz. They are the positive
Fourier coefficients of $|d|^2$. Applying
\eqref{I:B:eq:curvature-boundary-modulus} therefore gives
\[
 (H_dS_+\overline d)_{i-1}
   =\frac1{C^2}\sum_{j=0}^{m-i}p_{j+i}\overline{p_j}
 \quad(1\le i\le m),
\]
and these coefficients vanish for $i>m$. Consequently,
\begin{equation}
 T=\frac1{C^4}\sum_{i=1}^{m}
       \left|\sum_{j=0}^{m-i}p_{j+i}\overline{p_j}\right|^2.
 \label{I:B:eq:curvature-autocorrelation-P}
\end{equation}
Parseval's identity applied to the trigonometric polynomial $|P|^2$
shows that
\[
 \Bnorm P_{H^4}^4
 =\left(\sum_{j=0}^{m}|p_j|^2\right)^2
   +2\sum_{i=1}^{m}
       \left|\sum_{j=0}^{m-i}p_{j+i}\overline{p_j}\right|^2.
\]
The squared norm of $P$ is one, proving
\eqref{I:B:eq:curvature-hardy-coefficients}. Substitute these identities into
Proposition~\ref{I:B:prop:lift-curvature} to obtain
\eqref{I:B:eq:curvature-hardy-finite} and the norm expansion.
Finally, the scalar Hardy construction has
$E_{k-1}=1/C$, and its word lift has at most $k$ states. It is therefore
an admissible competitor for $\delta_k$, proving the upper bound in
\eqref{I:B:eq:curvature-hardy-norm}.
\end{proof}
